% Irodov Part IV §4.4 189--200. Source PDF 190--191, printed 181--182,
% answers PDF 332--333. Source figure 4.37/4.38.
\emwaveproblem{189}{181}
频率 $\nu=3.0\,\mathrm{MHz}$ 的电磁波从真空垂直进入非磁性介质，介电常数 $\varepsilon_r=4.0$。求波长增量。
\exsolution
真空中 $\lambda_0=c/\nu$，非磁介质相速度 $v=c/\sqrt{\varepsilon_r}$，
介质中波长 $\lambda=v/\nu$，故
\[
\boxed{\Delta\lambda=\lambda-\lambda_0
=\frac c\nu\left(\frac1{\sqrt{\varepsilon_r}}-1\right)}.
\]
代 $c=3.0\times10^8$m/s，
$\lambda_0=100$m、$\lambda=50$m，
\[
\boxed{\Delta\lambda=-50\,\mathrm m}.
\]
\excheck 频率越过静止界面时不变；变化的是相速度和波长。

\emwaveproblem{190}{181}
平面波垂直入射厚度为 $\ell$ 的平行板非磁介质，其相对介电常数从前表面 $\varepsilon_1$ 到后表面 $\varepsilon_2$ 按位置指数衰减。求某一固定波面穿越该板所需时间。
\exsolution
取 $x=0$ 为前表面，则
$\varepsilon(x)=\varepsilon_1
\exp[(x/\ell)\ln(\varepsilon_2/\varepsilon_1)]$。
局部相速度 $v(x)=c/\sqrt{\varepsilon(x)}$。
因此
\[
t=\frac1c\int_0^\ell\sqrt{\varepsilon(x)}\,\mathrm dx
=\frac{\ell\sqrt{\varepsilon_1}}c
\int_0^1 e^{(u/2)\ln(\varepsilon_2/\varepsilon_1)}\,\mathrm du.
\]
积分得到
\[
\boxed{t=
\frac{2\ell(\sqrt{\varepsilon_1}-\sqrt{\varepsilon_2})}
{c\ln(\varepsilon_1/\varepsilon_2)}}.
\]
\excheck $\varepsilon_2\to\varepsilon_1=\varepsilon$ 时取极限
$t=\ell\sqrt\varepsilon/c$，与均匀介质传播时间一致。

\emwaveproblem{191}{181}
频率 $\nu=10\,\mathrm{MHz}$ 的平面电磁波穿过弱导电介质，
电导率 $\sigma=10\,\mathrm{mS/m}$、相对介电常数 $\varepsilon_r=9$。求传导与位移电流密度振幅之比。
\exsolution
线性介质中传导电流
$\bm j_c=\sigma\bm E$，
位移电流
$\bm j_d=\partial_t\bm D
=\varepsilon_0\varepsilon_r\partial_t\bm E$。
若 $\bm E=\bm E_m\cos\omega t$，
振幅
$j_{c,m}=\sigma E_m$、
$j_{d,m}=\varepsilon_0\varepsilon_r\omega E_m$，
故
\[
\boxed{\frac{j_{c,m}}{j_{d,m}}
=\frac{\sigma}{2\pi\nu\varepsilon_0\varepsilon_r}
\simeq2.0}.
\]
\excheck 此比与电场幅度无关；两个电流分量相位相差 $\pi/2$，不能将它们的瞬时模直接相加。

\emwaveproblem{192}{181}
已知真空中平面电磁波的电场
$\bm E(\bm r,t)=\bm E_m\cos(\omega t-\bm k\cdot\bm r)$。
求 $\bm r=0$ 处的磁场强度矢量 $\bm H(t)$，
假设 $\bm E_m,\bm k$ 已知。
\exsolution
以法拉第方程
$\nabla\times\bm E=-\partial_t\bm B$：
$\nabla\times\bm E=
\bm k\times\bm E_m\sin(\omega t-\bm k\cdot\bm r)$，
故可取
$\bm B_m=(\bm k\times\bm E_m)/\omega$。
真空中 $\bm H=\bm B/\mu_0$，
$\omega=c|\bm k|$，于是
\[
\boxed{\bm H(0,t)=
\frac{\bm k\times\bm E_m}
{\mu_0c|\bm k|}\cos(\omega t)
=\sqrt{\frac{\varepsilon_0}{\mu_0}}
(\hat{\bm k}\times\bm E_m)\cos(\omega t)}.
\]
\excheck $\bm E$、$\bm H$、$\bm k$ 两两正交，
$\bm E\times\bm H$ 指向能量传播方向。

\emwaveproblem{193}{181}
真空平面电磁波电场
$\bm E=E_m\bm e_y\cos(\omega t-kx)$，
沿 $+x$ 传播。求位于
$\bm r=x\bm e_x$ 的磁场强度在
(a) $t=0$；(b) $t=t_0$。
取 $E_m=160\,\mathrm{V/m}$、
$k=0.51\,\mathrm{m^{-1}}$、
$x=7.7$m、$t_0=33$ns。
\exsolution
由第192题
\[
\boxed{\bm H(x,t)=\bm e_z E_m
\sqrt{\frac{\varepsilon_0}{\mu_0}}
\cos(ckt-kx)}.
\]
真空波阻抗
$Z_0=\sqrt{\mu_0/\varepsilon_0}\simeq377\,\Omega$，
故磁场强度幅度 $H_m=160/377\simeq0.424$A/m。
(a) 相位 $-kx=-3.927$：
$\boxed{\bm H(x,0)\simeq-0.30\bm e_z\,\mathrm{A/m}}$。
(b) $ckt_0-kx=(3.0\times10^8)(0.51)(33\times10^{-9})-3.927\simeq1.12$：
$\boxed{\bm H(x,t_0)\simeq+0.18\bm e_z\,\mathrm{A/m}}$。
\excheck 真空中给的是磁场强度 $\bm H$（A/m），若要磁感应强度须再乘 $\mu_0$。

\emwaveproblem{194}{181}
真空平面电磁波沿 $+x$ 传播，其电场
$\bm E=E_m\bm e_y\cos(\omega t-kx)$。
图4.37中边长 $\ell$ 的正方形导线框置于 $xy$ 平面，磁场沿 $z$。
求感应电动势的峰值；取
$E_m=0.50\,\mathrm{mV/m}$、
$\nu=5.0\,\mathrm{MHz}$、
$\ell=0.50$m。
\exsolution
磁感应强度
$B_z=(E_m/c)\cos(\omega t-kx)$，
$\omega=2\pi\nu$、$k=\omega/c$。
磁通
$\Phi(t)=\ell(E_m/c)\int_0^\ell\cos(\omega t-kx)\,\mathrm dx$，
因此
\[
\mathcal E=-\dot\Phi
=-\frac{\omega\ell E_m}{ck}
\left[\cos\omega t-\cos(\omega t-k\ell)\right].
\]
峰值的精确式是
\[
\boxed{\mathcal E_m=
2E_m\ell\left|\sin\frac{k\ell}{2}\right|}.
\]
由于 $k\ell\ll1$，近似为
\[
\boxed{\mathcal E_m\simeq\frac{\omega E_m\ell^2}{c}
=\frac{2\pi\nu E_m\ell^2}{c}
\simeq1.3\times10^{-5}\,\mathrm V=13\,\mathrm{\mu V}}.
\]
\excheck 源题明确写 $E_m=0.50\,\mathrm{mV/m}$，按 SI 单位计算应为 $13\,\mathrm{\mu V}$；书末印 $13\,\mathrm{mV}$，相差1000倍，需列为原书数值疑点。

\emwaveproblem{195}{182}
根据麦克斯韦方程证明图4.38所示真空平面波
$\bm E=E_y(x,t)\bm e_y$、
$\bm B=B_z(x,t)\bm e_z$
满足
$\partial_tE_y=-c^2\partial_xB_z$ 以及
$\partial_tB_z=-\partial_xE_y$。
\exsolution
法拉第方程
$\nabla\times\bm E=-\partial_t\bm B$。
在给定几何中旋度仅有 $z$ 分量
$(\nabla\times\bm E)_z=\partial_xE_y$，
所以
\[
\boxed{\partial_tB_z=-\partial_xE_y}.
\]
真空安培—麦克斯韦方程
$\nabla\times\bm B=\mu_0\varepsilon_0\partial_t\bm E$。
磁场只有 $B_z$ 分量，
其旋度 $y$ 分量为
$(\nabla\times\bm B)_y=-\partial_xB_z$，
故
$\partial_tE_y=-(\mu_0\varepsilon_0)^{-1}\partial_xB_z$，即
\[
\boxed{\partial_tE_y=-c^2\partial_xB_z}.
\]
\excheck 两式互相求导可分别得到
$\partial_t^2E=c^2\partial_x^2E$ 与
$\partial_t^2B=c^2\partial_x^2B$。

\emwaveproblem{196}{182}
真空平面电磁波
$\bm E=\bm E_m\cos(\omega t-\bm k\cdot\bm r)$。
求一周期平均坡印廷矢量。
\exsolution
根据平面波关系
$\bm H=(1/Z_0)\hat{\bm k}\times\bm E$，
能流密度
$\bm S=\bm E\times\bm H
=(E^2/Z_0)\hat{\bm k}$。
对周期取平均
$\langle\cos^2\rangle=1/2$：
\[
\boxed{\overline{\bm S}=
\frac12\varepsilon_0cE_m^2\hat{\bm k}
=\frac{\varepsilon_0c^2E_m^2}{2\omega}\bm k}.
\]
\excheck 单位为 W/m²，方向沿波矢；$E_m$ 是峰值不是有效值。

\emwaveproblem{197}{182}
真空中频率 $\nu=100\,\mathrm{MHz}$ 的线偏振平面电磁波，其电场振幅 $E_m=50\,\mathrm{mV/m}$。求：(a) 位移电流密度有效值；(b) 周期平均能流密度。
\exsolution
位移电流
$\bm j_d=\varepsilon_0\partial_t\bm E$，
幅度
$j_m=\varepsilon_0\omega E_m$。
正弦有效值为幅度除 $\sqrt2$，故
\[
\boxed{j_{d,\mathrm{rms}}=
\frac{\varepsilon_0(2\pi\nu)E_m}{\sqrt2}
\simeq0.20\,\mathrm{mA/m^2}}.
\]
平面波平均坡印廷通量
\[
\boxed{\overline S=
\frac12\varepsilon_0cE_m^2
\simeq3.3\,\mathrm{\mu W/m^2}}.
\]
\excheck $\varepsilon_0(2\pi10^8)(0.050)\simeq0.278$mA/m² 是\emph{振幅}，有效值再乘 $1/\sqrt2$。

\emwaveproblem{198}{182}
相对介电常数 $\varepsilon_r=4.0$ 的非磁介质中，振幅 $E_m=200$V/m 的平面电磁波掠过半径 $R=50$cm 的球，求 $t=1.0$min 内入射到球的总电磁能。
\exsolution
介质相速度 $v=c/\sqrt{\varepsilon_r}$，
波阻抗 $Z=Z_0/\sqrt{\varepsilon_r}$，
平均能流强度
$\overline S=E_m^2/(2Z)
=\tfrac12\varepsilon_0c\sqrt{\varepsilon_r}E_m^2$。
球在行波中的有效迎波面积为 $\pi R^2$，
故
\[
\boxed{W=\tfrac12\varepsilon_0c\sqrt{\varepsilon_r}\,
E_m^2\,\pi R^2t}.
\]
取 $t=60$s、$R=0.50$m，
$\boxed{W\simeq5.0\times10^3\,\mathrm J=5\,\mathrm{kJ}}$。
\excheck 不能以整个球表面积 $4\pi R^2$ 乘平面波能流强度；需使用迎波投影面积 $\pi R^2$。

\emwaveproblem{199}{182}
真空中沿 $x$ 建立电场驻波
$\bm E(x,t)=\bm e_yE_m\cos kx\cos\omega t$，
求磁感应强度的时空表达式，并描述
$t=0$、$t=T/4$ 时电磁场分布。
\exsolution
法拉第方程
$\partial_tB_z=-\partial_xE_y
=kE_m\sin kx\cos\omega t$。
积分得到振荡磁场
\[
\boxed{\bm B(x,t)=\bm e_z\frac{E_m}{c}
\sin kx\sin\omega t},
\quad \omega=ck.
\]
这选择无额外静磁场的纯驻波。
$t=0$ 时
$\bm E=\bm e_yE_m\cos kx$、$\bm B=0$；
$t=T/4$ 时 $\bm E=0$、
$\bm B=\bm e_z(E_m/c)\sin kx$。
\excheck 电、磁驻波的时间相位差四分之一周期，空间结点相互错开四分之一波长。

\emwaveproblem{200}{182}
真空电场驻波
$\bm E=E_m\bm e_y\cos kx\cos\omega t$。
求坡印廷矢量沿 x 的分量 $S_x(x,t)$，以及一周期的时间平均。
\exsolution
由第199题
$\bm H=\bm e_z(E_m/(\mu_0c))\sin kx\sin\omega t$。
坡印廷矢量
$\bm S=\bm E\times\bm H$，
因此
\[
S_x=
\frac{E_m^2}{\mu_0c}
\cos kx\sin kx\cos\omega t\sin\omega t.
\]
由积化和差并用 $(\mu_0c)^{-1}=\varepsilon_0c$：
\[
\boxed{S_x(x,t)=
\frac14\varepsilon_0cE_m^2\sin(2kx)\sin(2\omega t)}.
\]
对一完整周期平均
\[
\boxed{\overline{S_x}=0}.
\]
\excheck 驻波电磁能在空间往复交换，没有净单向远距离能流，但瞬时坡印廷矢量可不为零。
